\section{Related work and a practical map}
\label{sec:related}

Semantic agreement between informal mathematics and formal statements is now an
explicit target across several lines of work.  Statement-level methods include
learned or LLM-based critics such as FormalAlign, CriticLean, and Beyond
Compilation \citep{lu2025formalalign,peng2025criticlean,zhang2026beyondcompilation};
reference-based equivalence checks such as BEq/BEq+
\citep{liu2025beq,poiroux2025reliable}; and compiler-grounded tests such as
Theorem Testing and ShadowBench/SA-Pass
\citep{kim2026theoremtesting,han2026shadowbench}.  These methods help answer
whether a generated formal statement agrees with a target when an appropriate
reference, critic, dependent theorem set, or family of shadow statements is
available.

Project-scale systems address additional questions that arise once the source
and formal development both evolve.  EconCSLib provides paper-facing Lean
interfaces, dependency graphs, validation reports, and a human review dashboard
for checking translations of research papers \citep{garg2026econcs}.  Lean
Atlas/Lean Compass identifies the project declarations whose semantics can
affect selected target theorems and presents them in an interactive review
environment \citep{yanahama2026leanatlas}.  ATLAS/AutoformBot lets users compare
informal textbook statements with their Lean formalizations, inspect dependency
graphs, and view automated faithfulness assessments \citep{rammal2026scale}.
LeanMarathon, FormaTheoria, and LeanArchitect address adversarial target review,
long-horizon source/provenance management, and synchronization between informal
and formal project descriptions \citep{zhang2026leanmarathon,nie2026formatheoria,zhu2026leanarchitect}.

\begin{table}[t]
\centering
\small
\caption{A practical map for newcomers.  The systems overlap; each row names a
question for which the cited work provides a useful starting point.}
\label{tab:tool-map}
\begin{tabularx}{\linewidth}{@{}p{0.26\linewidth}p{0.34\linewidth}Y@{}}
\toprule
Question & Representative work & Main aid \\
\midrule
Does this formal statement match the intended theorem? & BEq/BEq+, CriticLean, ShadowBench & equivalence tests, semantic critics, shadow statements \\
What should a human inspect? & Lean Atlas / Compass & semantic review cone and interactive dependency view \\
How can an author review a paper translation? & EconCSLib & paper-to-Lean dashboard, LLM aid, human judgments \\
How can a large source corpus be browsed and evaluated? & ATLAS / AutoformBot & statement comparison, dependencies, faithfulness reports \\
How should review fit into a long project? & LeanMarathon, FormaTheoria, LeanArchitect & target review, provenance, reconciliation, blueprint maintenance \\
\bottomrule
\end{tabularx}
\end{table}

The local tooling in this paper grew independently from one Davis--Kahan
formalization and overlaps with several of these systems.  We use the comparison to organize practical lessons.  The paper makes no priority
claim for semantic review or review dashboards.
